\documentclass[../../main.tex]{subfiles}
\graphicspath{{../../image/}{../../figure/}}
\begin{document}

\section{初值在 $L^1$ 中}\label{section:ch12-L1估计}

回忆，由定理 11.5.2，
$L^1$ 中的初值诱导真非线性标量守恒律的解当 $t\to\infty$ 时以 $O(t^{-\frac12})$ 的速率衰减。此处旨在为真非线性两守恒律组的解建立类似性质。据此，我们考虑具有小振荡\eqref{align:ch12-解小振荡}、初值总变差不受限且属于 $L^1(-\infty,\infty)$ 的解 $(z(x,t),w(x,t))$：
\begin{align}
L=\int_{-\infty}^{\infty}[|z(x,0)|+|w(x,0)|]\,dx<\infty.
\label{align:ch12-L1}
\end{align}
主要结果是

\begin{theorem}\label{theorem:ch12-L1衰减}
当 $t\to\infty$ 时，
\begin{align}
(z(x,t),w(x,t))=O(t^{-\frac12}),
\label{align:ch12-L1衰减}
\end{align}
关于 $x\in(-\infty,\infty)$ 一致。
\end{theorem}

证明将分成若干步。

\begin{lemma}\label{lemma:ch12-L1积分}
对任意 $\bar t\in[0,\infty)$ 及 $-\infty<\bar x<\bar y<\infty$，
\begin{align}
\int_{\bar x}^{\bar y}[|z(x,\bar t)|+|w(x,\bar t)|]\,dx\leqslant 4\int_{\bar x-c\bar t}^{\bar y+c\bar t}[|z(x,0)|+|w(x,0)|]\,dx.
\label{align:ch12-L1积分}
\end{align}
特别地，$(z(\cdot,\bar t),w(\cdot,\bar t))\in L^1(-\infty,\infty)$。
\end{lemma}

\begin{proof}
我们通过求解\eqref{align:ch12-速矢}的 Goursat 问题并给定数据，构造一个 Lipschitz 连续熵 $\eta$：
\begin{align}
\begin{cases}
\eta(z,0)=|z|+\alpha z^2,\qquad -\infty<z<\infty,\\
\eta(0,w)=|w|+\alpha z^2,\qquad -\infty<w<\infty,
\end{cases}
\label{align:ch12-熵数据}
\end{align}
其中 $\alpha$ 是正常数。由\eqref{align:ch12-凸性}，当 $\alpha$ 充分大时，$\eta$ 是 $U$ 在含解值域的原点某邻域上的凸函数。

结合\eqref{align:ch12-速矢}与\eqref{align:ch12-熵数据}，当 $\delta$ 小时易推得
\begin{align}
\frac12(|z|+|w|)\leqslant\eta(z,w)\leqslant 2(|z|+|w|),\qquad -2\delta<z<2\delta,\ -2\delta<w<2\delta.
\label{align:ch12-熵等价}
\end{align}
进而，若 $q$ 是与 $\eta$ 相联系的熵流，归一化取 $q(0,0)=0$，则\eqref{align:ch12-熵数据}与\eqref{align:ch12-熵等价}蕴含
\begin{align}
|q(z,w)|\leqslant c\eta(z,w),\qquad -2\delta<z<2\delta,\ -2\delta<w<2\delta.
\label{align:ch12-熵流界}
\end{align}

现固定 $\bar t>0$、$-\infty<\bar x<\bar y<\infty$，将上述构造的熵--熵流对 $(\eta,q)$ 的不等式\eqref{align:ch12-熵不等式}在由 $\{(x,t):0<t<\bar t,\ \bar x-c(\bar t-t)<x<\bar y+c(\bar t-t)\}$ 定义的梯形区域上积分。利用\eqref{align:ch12-熵流界}，得
\begin{align}
\int_{\bar x}^{\bar y}\eta(z(x,\bar t),w(x,\bar t))\,dx\leqslant\int_{\bar x-c\bar t}^{\bar y+c\bar t}\eta(z(x,0),w(x,0))\,dx.
\label{align:ch12-梯形积分}
\end{align}
由\eqref{align:ch12-熵等价}，\eqref{align:ch12-梯形积分}蕴含\eqref{align:ch12-L1积分}。证明完成。
\end{proof}

\begin{lemma}\label{lemma:ch12-特征积分}
设 $(\bar z(\cdot),\bar w(\cdot))$ 表示沿上半平面任一点 $(\bar y,\bar t)$ 发出的极小（或极大）后向 1-（或 2-）特征线 $\xi(\cdot)$（或 $\zeta(\cdot)$）的 $(z,w)$ 迹。则
\begin{align}
\int_0^{\bar t}[\bar z^2(t)+|\bar w(t)|]\,dt\leqslant\tilde c L,
\label{align:ch12-特征积分1}\\
\int_0^{\bar t}[|\bar z(t)|+\bar w^2(t)]\,dt\leqslant\tilde c L.
\label{align:ch12-特征积分2}
\end{align}
\end{lemma}

\begin{proof}
只需验证\eqref{align:ch12-特征积分1}。设 $\eta$ 是任一 Lipschitz 连续凸熵，相联系熵流为 $q$，归一化取 $\eta(0,0)=0$、$q(0,0)=0$。我们固定 $\bar x<\bar y$，将不等式\eqref{align:ch12-熵不等式}在区域 $\{(x,t):0<t<\bar t,\ \bar x<x<\xi(t)\}$ 上积分，得
\begin{align}
\int_{\bar x}^{\bar y}\eta(z(x,\bar t),w(x,\bar t))\,dx-\int_{\bar x}^{\xi(0)}\eta(z(x,0),w(x,0))\,dx
+\int_0^{\bar t}G(\bar z(t),\bar w(t))\,dt-\int_0^{\bar t}q(z(\bar x+,t),w(\bar x+,t))\,dt\leqslant 0,
\label{align:ch12-区域积分-s08}
\end{align}
其中 $G$ 定义为
\begin{align}
G(z,w)=q(z,w)-\lambda(z,w)\eta(z,w).
\label{align:ch12-G定义}
\end{align}

我们寻找一个使 $G(z,w)$ 在 $(-2\delta,2\delta)\times(-2\delta,2\delta)$ 上正定的熵--熵流对。由\eqref{align:ch12-熵数据}，
\begin{align}
G_z&=-\lambda_z\eta,
\label{align:ch12-Gz}\\
G_w&=[(\mu-\lambda)\eta]_w-\mu_w\eta,
\label{align:ch12-Gw}
\end{align}
这表明当 $z\to0$ 时 $G$ 至少二次衰减，但当 $w\to0$ 时可能衰减更慢、甚至线性。

我们通过求解\eqref{align:ch12-速矢}的 Goursat 问题并给定数据构造熵 $\eta$：
\begin{align}
\begin{cases}
\eta(z,0)=2z+\alpha z^2,\qquad -\infty<z<\infty,\\
\eta(0,w)=|w|+\alpha w^2,\qquad -\infty<w<\infty.
\end{cases}
\label{align:ch12-熵数据2}
\end{align}
当 $\alpha$ 充分大时，由\eqref{align:ch12-凸性}，$\eta$ 是 $U$ 在含解值域的原点某邻域上的凸函数。由\eqref{align:ch12-Gw}、\eqref{align:ch12-速矢}与\eqref{align:ch12-熵数据2}，推得
\begin{align}
G(0,w)&=[\mu(0,0)-\lambda(0,0)]|w|+O(w^2),
\label{align:ch12-G0w}\\
\eta(z,w)&=2z+|w|+O(z^2+w^2),
\label{align:ch12-eta展开}
\end{align}
对 $(z,w)$ 靠近原点成立。结合\eqref{align:ch12-G0w}与\eqref{align:ch12-Gz}、\eqref{align:ch12-eta展开}，我们断定
\begin{align}
G(z,w)=[\mu(0,0)-\lambda(0,0)]|w|-\lambda_z(0,0)z^2+O(w^2+|zw|+|z|^3).
\label{align:ch12-G展开}
\end{align}

现回到\eqref{align:ch12-区域积分}。由\cref{lemma:ch12-L1积分}，对所有 $t\in[0,\bar t]$，$(z(\cdot,t),w(\cdot,t))\in L^1(-\infty,\infty)$，故
\begin{align}
\liminf_{\bar x\to-\infty}\left|\int_0^{\bar t}q(z(\bar x+,t),w(\bar x+,t))\,dt\right|=0.
\label{align:ch12-q极限}
\end{align}
因此，\eqref{align:ch12-区域积分}、\eqref{align:ch12-q极限}、\eqref{align:ch12-eta展开}、\eqref{align:ch12-L1积分}与\eqref{align:ch12-L1}共同蕴含
\begin{align}
\int_0^{\bar t}G(\bar z(t),\bar w(t))\,dt\leqslant 12L,
\label{align:ch12-G积分}
\end{align}
只要\eqref{align:ch12-解小振荡} 以充分小的 $\delta$ 成立。由\eqref{align:ch12-G积分}、\eqref{align:ch12-G展开}与\eqref{align:ch12-速度差}，断言\eqref{align:ch12-特征积分1} 现在容易推出。证明完成。
\end{proof}

\cref{lemma:ch12-特征积分} 表明：沿极小后向 1-特征线，$z$ 为 $O(t^{-\frac12})$、$w$ 为 $O(t^{-1})$；而沿极大后向 2-特征线，$z$ 为 $O(t^{-1})$、$w$ 为 $O(t^{-\frac12})$。事实上，回忆 $\bar z(\cdot)$ 与 $\bar w(\cdot)$ 分别沿极小与极大后向 1-、2-特征线非增，我们直接从\eqref{align:ch12-特征积分1}与\eqref{align:ch12-特征积分2}推知，当 $t\to\infty$ 时 $z(x,t)$ 与 $w(x,t)$ 的正部为 $O(t^{-\frac12})$。现通过建立两侧的 $O(t^{-\frac12})$ 衰减来完成\cref{theorem:ch12-L1衰减}的证明：

\begin{lemma}\label{lemma:ch12-逐点衰减}
当 $\delta$ 充分小时，
\begin{align}
z^2(x,t)&\leqslant\frac{8\tilde c L}{t},
\label{align:ch12-逐点衰减1}\\
w^2(x,t)&\leqslant\frac{8\tilde c L}{t},
\label{align:ch12-逐点衰减2}
\end{align}
对所有 $-\infty<x<\infty$、$0<t<\infty$ 成立，其中 $\tilde c$ 是\eqref{align:ch12-特征积分1}与\eqref{align:ch12-特征积分2}中的常数。
\end{lemma}

\begin{proof}
反证，设断言不真，令 $\bar t>0$ 为使\eqref{align:ch12-逐点衰减1}和/或\eqref{align:ch12-逐点衰减2}对某 $x$ 被破坏的点集的最大下界。由\cref{theorem:ch12-局部正则存在}，解在 $\bar t$ 之外的延拓是通过沿 $\bar t$-时间线求解 Riemann 问题而启动的。因此，由于\eqref{align:ch12-逐点衰减1}和/或\eqref{align:ch12-逐点衰减2}对 $t>\bar t$ 失效，可找到 $\bar y\in(-\infty,\infty)$ 使
\begin{align}
z^2(\bar y,\bar t)&>\frac{4\tilde c L}{\bar t},
\label{align:ch12-破坏1}\\
w^2(\bar y,\bar t)&>\frac{4\tilde c L}{\bar t},
\label{align:ch12-破坏2}
\end{align}
和/或成立。为确定起见，设\eqref{align:ch12-破坏1} 成立。

令 $(\bar z(\cdot),\bar w(\cdot))$ 表示沿从 $(\bar y,\bar t)$ 发出的极小后向 1-特征线 $\xi(\cdot)$ 的 $(z,w)$ 迹。将\cref{theorem:ch12-TV界}用于时间原点由 $t=0$ 平移到 $t=\bar t/2$ 的情形，推得
\begin{align}
TV_{[\frac12\bar t,\bar t]}\bar w(\cdot)\leqslant\tilde c\{TV_{[\bar y-\frac12\bar\mu\bar t,\bar y-\frac12\bar\lambda\bar t]}z\left(\cdot,\frac{\bar t}{2}\right)+TV_{[\bar y-\frac12\bar\mu\bar t,\bar y-\frac12\bar\lambda\bar t]}w\left(\cdot,\frac{\bar t}{2}\right)\},
\label{align:ch12-TVw}
\end{align}
其中 $\bar\lambda$ 是解值域上 $\lambda(z,w)$ 的下确界，$\bar\mu$ 是 $\mu(z,w)$ 的上确界。我们借助\cref{theorem:ch12-TV散开}估计\eqref{align:ch12-TVw}右端，得
\begin{align}
TV_{[\frac12\bar t,\bar t]}\bar w(\cdot)\leqslant\tilde c[b(\bar\mu-\bar\lambda)+\beta\delta].
\label{align:ch12-TVw估计}
\end{align}
由假设，
\begin{align}
\bar w^2(t)\leqslant\frac{16\tilde c L}{\bar t},\qquad \frac{\bar t}{2}\leqslant t<\bar t.
\label{align:ch12-w界-s08}
\end{align}
我们还有 $|\bar z(t)|\leqslant 2\delta$。因此，应用\eqref{align:ch12-z等式}推得
\begin{align}
\bar z^2(\bar t-)-\bar z^2(t)\leqslant\bar c\delta\frac{4\tilde c L}{\bar t},
\label{align:ch12-z平方}
\end{align}
其中 $\bar c=64a\hat c[b(\bar\mu-\bar\lambda)+\beta\delta]$。

由于 $\bar z(\bar t-)=z(\bar y,\bar t)$，结合\eqref{align:ch12-破坏1}与\eqref{align:ch12-z平方}给出
\begin{align}
\bar z^2(t)\geqslant\frac{4\tilde c L}{\bar t}(1-\bar c\delta),\qquad \frac{\bar t}{2}\leqslant t<\bar t.
\label{align:ch12-z下界}
\end{align}
由\eqref{align:ch12-z下界}，
\begin{align}
\int_{\frac{\bar t}{2}}^{\bar t}\bar z^2(t)\,dt\geqslant 2\tilde c L(1-\bar c\delta),
\label{align:ch12-z积分}
\end{align}
当 $\delta$ 充分小时，这与\eqref{align:ch12-特征积分1}构成期望的矛盾。证明完成。
\end{proof}

\end{document}
